% GENERATED FILE. Edit canonical lessons, then run npm run generate:books.
% canonical-source-hash: fnv1a64:5c0cb97c7b6a0c54
% canonical-lessons: ZH-C16-hear-zi, ZH-C16-zi, ZH-C16-hear-hanzi, ZH-W16-han-observe, ZH-W16-han-guided, ZH-W16-han-delayed, ZH-C16-hanzi, ZH-C16-hear-hanyu, ZH-R16-identity-1, ZH-W16-yu-observe, ZH-W16-yu-guided, ZH-W16-yu-delayed, ZH-C16-hanyu, ZH-R16-identity-2, ZH-C16-practice

\chapter{Character or language?}
\label{ch:zh-language-identity}
\hlchaptermodality{16}

\begin{chapteropening}
\textbf{By the end of this chapter:} \emph{I can understand, say, read and write \zh{字}, \zh{汉字} and \zh{汉语} after gaining \zh{汉} and \zh{语} through source-verified observation, tracing, guided copying and delayed recall.}

Three identity words are retrieved across two expanding reviews, then checked independently in listening, speaking, unassisted reading and no-model writing.
\end{chapteropening}

\section[zì]{Hear \emph{zì} — a written character}

\label{lesson:ZH-C16-hear-zi}



\begin{quote}
Say \textbf{middle school} in Mandarin.
\end{quote}

\subsection*{What to know first}
\begin{quote}
\emph{zì} — \textbf{a written Chinese character}
\end{quote}

You heard a light \emph{zi} at the end of \emph{míngzi}. Standing alone, this word has a sharp falling tone: \emph{zì}. Stay with sound and meaning for this tiny step.

\begin{tcolorbox}[breakable,colback=teal!4,colframe=teal!35!black,title={Before you move on}]
Source: \href{https://hsk.cn-bj.ufileos.com/3.0/\%E6\%96\%B0\%E7\%89\%88HSK\%E8\%80\%83\%E8\%AF\%95\%E5\%A4\%A7\%E7\%BA\%B21219.pdf}{Official HSK 3.0 syllabus, Level 1 vocabulary}.
\end{tcolorbox}
\section[zì]{\zh{字} — the familiar shape stands alone}

\label{lesson:ZH-C16-zi}



\begin{quote}
Say \textbf{written character} in Mandarin.
\end{quote}

\begin{cousinweb}
\begin{quote}
\textbf{\zh{字}} — \emph{zì} — \textbf{a written Chinese character}
\end{quote}

Hide the pinyin. Read \zh{字} from its familiar roof-over-child shape, then say its falling tone. You are attaching a standalone word to writing you already know.
\end{cousinweb}

\subsection*{Your turn}
Look at \textbf{\zh{字}} for three seconds. Cover it, wait five seconds, then write it once. Check only two things: roof before child, and six strokes.

\begin{tcolorbox}[breakable,colback=teal!4,colframe=teal!35!black,title={Before you move on}]
Source: \href{https://hsk.cn-bj.ufileos.com/3.0/\%E6\%96\%B0\%E7\%89\%88HSK\%E8\%80\%83\%E8\%AF\%95\%E5\%A4\%A7\%E7\%BA\%B21219.pdf}{Official HSK 3.0 syllabus, Level 1 vocabulary}.
\end{tcolorbox}
\section[hànzì]{Hear \emph{hànzì} — Chinese character}

\label{lesson:ZH-C16-hear-hanzi}



\begin{quote}
Say \emph{zì} and its meaning.
\end{quote}

\subsection*{What to know first}
\begin{quote}
\emph{hànzì} — \textbf{Chinese character; Han character}
\end{quote}

The second half is your known \emph{zì}, written character. The first half, \emph{hàn}, labels it as Han Chinese. Hear two clean falling tones: \emph{hàn-zì}. The new first character stays hidden until the next lesson.

\begin{tcolorbox}[breakable,colback=teal!4,colframe=teal!35!black,title={Before you move on}]
Source: \href{https://hsk.cn-bj.ufileos.com/3.0/\%E6\%96\%B0\%E7\%89\%88HSK\%E8\%80\%83\%E8\%AF\%95\%E5\%A4\%A7\%E7\%BA\%B21219.pdf}{Official HSK 3.0 syllabus, Level 1 vocabulary}.
\end{tcolorbox}
\section[hàn]{\zh{汉} — watch five strokes}

\label{lesson:ZH-W16-han-observe}



\begin{quote}
Say \textbf{Chinese character}, then point to the new first half of \zh{汉字}.
\end{quote}

\subsection*{Script — the shape}
In \textbf{\zh{汉}}, write \zh{氵} first. On the right, stroke four crosses right and turns down-left without lifting. Stroke five falls down-right. Five strokes; four lifts. Watch the crossing open into two different falling directions.

\subsection*{Your turn}
Trace \textbf{\zh{汉}} once, slowly. Say \textbf{dot, dot, rise; turn left; fall right}.

\begin{tcolorbox}[breakable,colback=teal!4,colframe=teal!35!black,title={Before you move on}]
Source: \href{https://raw.githubusercontent.com/chanind/hanzi-writer-data/68d10a4b21150cae5e1ebbd223eed289cf32d90c/data/\%E6\%B1\%89.json}{Hanzi Writer Data \zh{汉}, pinned stroke paths and medians}.
\end{tcolorbox}
\section[hàn]{\zh{汉} — copy with a five-part cue}

\label{lesson:ZH-W16-han-guided}



\begin{quote}
Air-write \zh{氵}: dot, dot, rise.
\end{quote}

\subsection*{Your turn}
Copy \textbf{\zh{汉}} twice while the model remains visible. First say the five-part cue. On the second copy, pause after \zh{氵} and check that the right side still has only two strokes.

\begin{tcolorbox}[breakable,colback=teal!4,colframe=teal!35!black,title={Before you move on}]
Source: \href{https://raw.githubusercontent.com/chanind/hanzi-writer-data/68d10a4b21150cae5e1ebbd223eed289cf32d90c/data/\%E6\%B1\%89.json}{Hanzi Writer Data \zh{汉}, pinned stroke paths and medians}.
\end{tcolorbox}
\section[hàn]{\zh{汉} — hold the route briefly}

\label{lesson:ZH-W16-han-delayed}



\begin{quote}
Look at \textbf{\zh{汉}} for five seconds. Cover it.
\end{quote}

\subsection*{Your turn}
Study \textbf{\zh{汉}} for five seconds. Cover it. Say the five-part cue, wait five more seconds, then write once. Reveal the model and repair only the first wrong stroke; do not copy the whole character again unless needed.

\begin{tcolorbox}[breakable,colback=teal!4,colframe=teal!35!black,title={Before you move on}]
Source: \href{https://raw.githubusercontent.com/chanind/hanzi-writer-data/68d10a4b21150cae5e1ebbd223eed289cf32d90c/data/\%E6\%B1\%89.json}{Hanzi Writer Data \zh{汉}, pinned stroke paths and medians}.
\end{tcolorbox}
\section[hànzì]{\zh{汉字} — two known routes, one word}

\label{lesson:ZH-C16-hanzi}



\begin{quote}
Say \textbf{Chinese character} in Mandarin.
\end{quote}

\begin{cousinweb}
\begin{quote}
\textbf{\zh{汉字}} — \emph{hànzì} — \textbf{Chinese character}
\end{quote}

Read the two characters first without pinyin: Han Chinese + written character.
\end{cousinweb}

\subsection*{Your turn}
Copy \textbf{\zh{汉字}} once. Leave a small, even gap; do not merge the two squares.

\begin{tcolorbox}[breakable,colback=teal!4,colframe=teal!35!black,title={Before you move on}]
Source: \href{https://hsk.cn-bj.ufileos.com/3.0/\%E6\%96\%B0\%E7\%89\%88HSK\%E8\%80\%83\%E8\%AF\%95\%E5\%A4\%A7\%E7\%BA\%B21219.pdf}{Official HSK 3.0 syllabus, Level 1 vocabulary}.
\end{tcolorbox}
\section[hànyǔ]{Hear \emph{hànyǔ} — the Chinese language}

\label{lesson:ZH-C16-hear-hanyu}



\begin{quote}
Say \textbf{Chinese character} in Mandarin.
\end{quote}

\subsection*{What to know first}
\begin{quote}
\emph{hànyǔ} — \textbf{the Chinese language; Mandarin}
\end{quote}

Keep the known falling \emph{hàn}. Then let \emph{yǔ} dip in third tone. The new written half waits; today the contrast is heard: \emph{hànzì} is a character, \emph{hànyǔ} is the language.

\begin{tcolorbox}[breakable,colback=teal!4,colframe=teal!35!black,title={Before you move on}]
Source: \href{https://hsk.cn-bj.ufileos.com/3.0/\%E6\%96\%B0\%E7\%89\%88HSK\%E8\%80\%83\%E8\%AF\%95\%E5\%A4\%A7\%E7\%BA\%B21219.pdf}{Official HSK 3.0 syllabus, Level 1 vocabulary}.
\end{tcolorbox}
\section[Practice]{Retrieval 1 — \zh{字} and \zh{汉字}}

\label{lesson:ZH-R16-identity-1}



\begin{quote}
Say and write \textbf{written character} without a model.
\end{quote}

\subsection*{Your turn}
Hear \textbf{written character} and \textbf{Chinese character} in mixed order. Say each answer, read its unpointed card, then write it without a model. Repair only the missed item.

\begin{tcolorbox}[breakable,colback=teal!4,colframe=teal!35!black,title={Before you move on}]
Source: \href{https://hsk.cn-bj.ufileos.com/3.0/\%E6\%96\%B0\%E7\%89\%88HSK\%E8\%80\%83\%E8\%AF\%95\%E5\%A4\%A7\%E7\%BA\%B21219.pdf}{Official HSK 3.0 syllabus, Level 1 vocabulary}.
\end{tcolorbox}
\section[yǔ]{\zh{语} — speech, five, box}

\label{lesson:ZH-W16-yu-observe}



\begin{quote}
Air-write \zh{讠}, then say \emph{hànyǔ}.
\end{quote}

\subsection*{Script — the shape}
Write \textbf{\zh{讠}} first. On the right, build \textbf{\zh{五}} from top to bottom, then close \textbf{\zh{口}} at the base. Nine strokes; eight lifts. Keep the speech side narrow so the right side has room.

\subsection*{Your turn}
Trace \textbf{\zh{语}} once. Say \textbf{speech, five, box} as your hand crosses the three parts.

\begin{tcolorbox}[breakable,colback=teal!4,colframe=teal!35!black,title={Before you move on}]
Source: \href{https://raw.githubusercontent.com/chanind/hanzi-writer-data/68d10a4b21150cae5e1ebbd223eed289cf32d90c/data/\%E8\%AF\%AD.json}{Hanzi Writer Data \zh{语}, pinned stroke paths and medians}.
\end{tcolorbox}
\section[yǔ]{\zh{语} — copy, then fade}

\label{lesson:ZH-W16-yu-guided}



\begin{quote}
Say \textbf{the Chinese language}, then point across \textbf{\zh{语}} and say \textbf{speech, five, box}.
\end{quote}

\subsection*{Your turn}
Copy \textbf{\zh{语}} once beside the model. Cover it halfway. Copy again from the visible half, keeping \zh{讠} narrow and the final \zh{口} closed.

\begin{tcolorbox}[breakable,colback=teal!4,colframe=teal!35!black,title={Before you move on}]
Source: \href{https://raw.githubusercontent.com/chanind/hanzi-writer-data/68d10a4b21150cae5e1ebbd223eed289cf32d90c/data/\%E8\%AF\%AD.json}{Hanzi Writer Data \zh{语}, pinned stroke paths and medians}.
\end{tcolorbox}
\section[yǔ]{\zh{语} — wait, then write}

\label{lesson:ZH-W16-yu-delayed}



\begin{quote}
Say \textbf{the Chinese language}. Look at \textbf{\zh{语}} for five seconds, then cover it.
\end{quote}

\subsection*{Your turn}
Tap once for every lift while saying \textbf{speech, five, box}. Now write \textbf{\zh{语}} without uncovering the model. Compare only after you finish.

\begin{tcolorbox}[breakable,colback=teal!4,colframe=teal!35!black,title={Before you move on}]
Source: \href{https://raw.githubusercontent.com/chanind/hanzi-writer-data/68d10a4b21150cae5e1ebbd223eed289cf32d90c/data/\%E8\%AF\%AD.json}{Hanzi Writer Data \zh{语}, pinned stroke paths and medians}.
\end{tcolorbox}
\section[hànyǔ]{\zh{汉语} — the language in writing}

\label{lesson:ZH-C16-hanyu}



\begin{quote}
Say \textbf{the Chinese language} in Mandarin.
\end{quote}

\begin{cousinweb}
Read this before looking below:

\begin{quote}
\textbf{\zh{汉语}}
\end{quote}

It joins Han Chinese + language: \emph{hànyǔ}, \textbf{the Chinese language}.
\end{cousinweb}

\subsection*{Your turn}
Copy \textbf{\zh{汉语}} once. Keep two separate squares and read the whole word aloud.

\begin{tcolorbox}[breakable,colback=teal!4,colframe=teal!35!black,title={Before you move on}]
Source: \href{https://hsk.cn-bj.ufileos.com/3.0/\%E6\%96\%B0\%E7\%89\%88HSK\%E8\%80\%83\%E8\%AF\%95\%E5\%A4\%A7\%E7\%BA\%B21219.pdf}{Official HSK 3.0 syllabus, Level 1 vocabulary}.
\end{tcolorbox}
\section[Practice]{Retrieval 2 — character or language?}

\label{lesson:ZH-R16-identity-2}



\begin{quote}
Say and write \textbf{the Chinese language} without a model.
\end{quote}

\subsection*{Your turn}
Hear \textbf{written character}, \textbf{Chinese character}, and \textbf{Chinese language} in a mixed order. Say, read, and write each answer. For the two \zh{汉} compounds, listen for the ending before choosing \zh{字} or \zh{语}.

\begin{tcolorbox}[breakable,colback=teal!4,colframe=teal!35!black,title={Before you move on}]
Source: \href{https://hsk.cn-bj.ufileos.com/3.0/\%E6\%96\%B0\%E7\%89\%88HSK\%E8\%80\%83\%E8\%AF\%95\%E5\%A4\%A7\%E7\%BA\%B21219.pdf}{Official HSK 3.0 syllabus, Level 1 vocabulary}.
\end{tcolorbox}
\section[Practice]{Checkpoint — \zh{字}, \zh{汉字}, \zh{汉语}}

\label{lesson:ZH-C16-practice}



\begin{quote}
Write \zh{汉} and \zh{语} without a model.
\end{quote}

\subsection*{Your turn}
1. \textbf{Listening:} hear the three words in mixed order and choose each meaning. 2. \textbf{Speaking:} produce all three from meaning cards, with tones audible. 3. \textbf{Reading:} read \textbf{\zh{字}, \zh{汉语}, \zh{汉字}} cold and give each meaning. 4. \textbf{Writing:} hear \textbf{Chinese character} and \textbf{Chinese language}, then write both without a model.

Pass each skill separately. Check the \zh{字} or \zh{语} ending only after both written words are complete.

\begin{tcolorbox}[breakable,colback=teal!4,colframe=teal!35!black,title={Before you move on}]
Source: \href{https://hsk.cn-bj.ufileos.com/3.0/\%E6\%96\%B0\%E7\%89\%88HSK\%E8\%80\%83\%E8\%AF\%95\%E5\%A4\%A7\%E7\%BA\%B21219.pdf}{Official HSK 3.0 syllabus, Level 1 vocabulary}.
\end{tcolorbox}
